两组深度图都能表现出建筑的前后关系，主体轮廓比较接近。例如 0025 的平均置信度从 0.846 升至 0.861，而 0043 从 0.757 降至 0.747，说明变化与观察位置有关。5 视图增加了参与匹配的图像，但并没有使各项统计都同时变好：平均置信度只从 0.734 略增至 0.736，超过 0.8 的像素比例反而从 68.13\% 降至 67.38\%。因此，\textbf{增加输入图像数量，并不意味着每个位置都更有把握。}

经过相同的几何过滤后，5 视图的最终保留比例为 64.25\%，略高于 3 视图的 63.63\%，两组融合点数则十分接近。更多视角可以补充对应关系，但也可能带来遮挡或外观差异。本次结果说明两组都能完成重建，5 视图在部分位置提供了帮助；没有真实参考值时，不能仅凭这几项内部统计认定它的绝对精度一定更高。
